%! Mean, Variance, and Covariance — Unit 8, Lesson 8.4 — Slides (Beamer)
\documentclass[aspectratio=169, 11pt]{beamer}
\usepackage{linalg-beamer}   % provides \burgundyheader{} and \sectionlabel[color]{}

\begin{document}

% TITLE SLIDE (CourseName is not defined in beamer, so the name is literal)
{
\setbeamercolor{background canvas}{bg=burgundy}
\begin{frame}[plain]
  \vspace{2.8cm}\hspace{0.55cm}%
  \begin{minipage}{0.9\textwidth}
    {\color{goldacc}\bfseries\small LESSON 8.4}\\[0.5cm]
    {\color{white}\bfseries\LARGE Mean, Variance, and Covariance}\\[1.0cm]
    {\color{blushmid}\small Linear Algebra~$\cdot$~Unit 8: Learning from Data}
  \end{minipage}
\end{frame}
}

% HOOK
\begin{frame}[t, plain]
  \burgundyheader{What is the shape of a data set?}
  \sectionlabel{Hook}
  \begin{columns}[T]
    \begin{column}{0.54\textwidth}
      Four students, two quiz scores each: $(2,4),(4,2),(6,8),(8,6)$ --- a \textbf{cloud of points}.\\[6pt]
      Describe its shape with a few numbers:
      \begin{itemize}
        \item \textbf{where} it centers --- the \emph{mean};
        \item \textbf{how spread} it is --- the \emph{variance};
        \item \textbf{which way it leans} --- the \emph{covariance}.
      \end{itemize}
    \end{column}
    \begin{column}{0.46\textwidth}
      \centering
      \begin{tikzpicture}[scale=0.62]
        \draw[linegray, dashed] (0.2,2)--(3.8,2);
        \draw[linegray, dashed] (2,0.2)--(2,3.8);
        \node[charcoal] at (3.6,2.0) {\scriptsize $x$};
        \node[charcoal] at (2.0,3.75) {\scriptsize $y$};
        \fill[royalblue] (0.8,1.6) circle (2.6pt);
        \fill[royalblue] (1.6,0.8) circle (2.6pt);
        \fill[royalblue] (2.4,3.2) circle (2.6pt);
        \fill[royalblue] (3.2,2.4) circle (2.6pt);
        \draw[->, very thick, burgundy] (2,2)--(3.2,3.2);
        \node[burgundy] at (3.0,3.45) {\scriptsize leans up-right};
      \end{tikzpicture}
    \end{column}
  \end{columns}
  \vspace{2pt}
  \begin{block}{The question of Lesson 8.4}
    Which few numbers capture where data sits, how it spreads, and how its features move together?
  \end{block}
\end{frame}

% MEAN + VARIANCE
\begin{frame}[t, plain]
  \burgundyheader{Mean $=$ center, \; Variance $=$ spread}
  \sectionlabel[goldacc]{Idea}
  \begin{columns}[T]
    \begin{column}{0.5\textwidth}
      \textbf{Mean} (the center):
      \[ \mu=\tfrac1N\textstyle\sum x_i. \]
      For $2,4,6,8$: $\mu=5$.\\[8pt]
      \textbf{Variance} (the spread) $=$ average squared distance from the mean:
      \[ \sigma^2=\tfrac1N\textstyle\sum (x_i-\mu)^2. \]
    \end{column}
    \begin{column}{0.5\textwidth}
      For $2,4,6,8$: deviations $-3,-1,1,3$,
      \[ \sigma^2=\frac{9+1+1+9}{4}=5,\quad \sigma=\sqrt5. \]
      \vspace{2pt}
      We \emph{square} the distances so $+$ and $-$ gaps don't cancel --- the same move as Unit-4 least squares.\\[6pt]
      Same mean, bigger variance $=$ more scattered.
    \end{column}
  \end{columns}
\end{frame}

% COVARIANCE + THE MATRIX
\begin{frame}[t, plain]
  \burgundyheader{Covariance: do two features move together?}
  \sectionlabel{Skill}
  \[ \sigma_{xy}=\tfrac1N\textstyle\sum (x_i-\mu_x)(y_i-\mu_y)
     \qquad \begin{cases}>0 & \text{together}\\ <0 & \text{opposite}\\ =0 & \text{unrelated}\end{cases} \]
  Spine $(2,4),(4,2),(6,8),(8,6)$: centered $x:-3,-1,1,3$, $y:-1,-3,3,1$, so
  $\sigma_{xy}=\tfrac{3+3+3+3}{4}=\mathbf{3}>0$. Pack it into the symmetric \textbf{covariance matrix}:
  \[ V=\begin{bmatrix}\sigma_x^2 & \sigma_{xy}\\ \sigma_{xy} & \sigma_y^2\end{bmatrix}
      =\begin{bmatrix}5 & 3\\ 3 & 5\end{bmatrix}. \]
  \vspace{-2pt}
  Variances on the diagonal, covariance off it --- and $V$ is \textbf{symmetric} ($\sigma_{xy}=\sigma_{yx}$).
\end{frame}

% THE FINALE
\begin{frame}[t, plain]
  \burgundyheader{The covariance matrix $=$ the whole course}
  \sectionlabel[goldacc]{Payoff}
  \begin{itemize}
    \setlength{\itemsep}{6pt}
    \item $V=\left[\begin{smallmatrix}5&3\\3&5\end{smallmatrix}\right]$ is \textbf{symmetric} (Unit~6):
          real eigenvalues $\lambda=8,2$; perpendicular eigenvectors $(1,1),(1,-1)$.
    \item Those eigenvectors are the \textbf{principal components} (Unit~7): PC1 $=(1,1)$ holds
          $\tfrac{8}{10}=80\%$ of the spread. Total variance $=$ trace $=5+5=10=8+2$.
    \item \textbf{The finish line:} raw \emph{data} (U8) $\to$ symmetric \emph{covariance matrix} (U6) $\to$
          \emph{eigen-directions} $=$ PCA (U7), all from means, distances, and dot products (U1, U4).
  \end{itemize}
  \vspace{2pt}
  \begin{block}{End of the course}
    Linear algebra is the language for the \textbf{shape of data}. \emph{Well done.}
  \end{block}
\end{frame}

\end{document}
